% ============================================================
%  CAPITOLO 5 --- PROGETTAZIONE DELL'APPLICAZIONE
% ============================================================
\chapter{Progettazione dell'applicazione}

\section{Architettura del sistema}

L'applicazione segue un'architettura \textbf{client-server a tre livelli}:

\begin{enumerate}
  \item \textbf{Presentation layer}: interfaccia web con interattivit\`{a} parziale tramite \textit{HTMX}
  \item \textbf{Application layer}: Applicazione Django
  \item \textbf{Data layer}: DBMS MySQL
\end{enumerate}

\section{Tecnologie utilizzate}

\begin{table}[H]
  \centering
  \caption{Stack tecnologico}
  \label{tab:stack}
  \rowcolors{2}{rowalt}{white}
  \begin{tabularx}{\textwidth}{L{3cm} L{3.5cm} L{8.2cm}}
    \toprule
    \rowcolor{primary!15}
    \textbf{Livello} & \textbf{Tecnologia} & \textbf{Ruolo} \\
    \midrule
    Database     & MySQL            & DBMS relazionale, gestione dati \\
    Connettore   & mysqlclient      & Driver di connessione Python $\rightarrow$ MySQL \\
    Backend      & Python 3.13      & Linguaggio applicativo principale \\
    Framework    & Django 6.0       & Routing HTTP, gestione sessioni, template engine \\
    Accesso dati & Query SQL pure   & Interrogazioni dirette al database, senza ORM \\
    Frontend     & HTML5 + CSS3     & Interfaccia utente \\
    Interattivit\`{a} & HTMX 1.9    & Aggiornamento parziale della pagina via richieste AJAX, senza framework JS pesanti \\
    Stile        & CSS custom       & Tema grafico personalizzato, nessun framework di terze parti \\
    \bottomrule
  \end{tabularx}
\end{table}

\clearpage

\section{Struttura del progetto}

\begin{verbatim}
appStabilimento/
|---- manage.py                      # Entry point Django
|---- pyproject.toml / uv.lock       # Gestione dipendenze (uv)
`---- gestionaleStabilimento/
    |---- settings.py                # Configurazione DB, static, sessioni
    |---- urls.py                    # Definizione di tutte le rotte
    |---- views/
    |   |---- auth.py                # Login/logout/dashboard
    |   |---- spiaggia.py            # Clienti, prenotazioni, abbonamenti, noleggi
    |   |---- ordini.py              # Fornitori, ordini, magazzino
    |   `---- turni.py               # Dipendenti e turni
    |---- templates/
    |   |---- base.html              # Layout comune, da estendere
    |   |---- login.html
    |   |---- dashboard.html
    |   |---- partials/
    |   |   `---- _toast_oob.html    # Notifiche riusabili (successo/errore)
    |   |---- spiaggia/
    |   |   |---- main.html
    |   |   `---- partials/          # Form e risultati query dell'area spiaggia
    |   |---- ordini/
    |   |   |---- main.html
    |   |   `---- partials/
    |   `---- turni/
    |       |---- main.html
    |       `---- partials/
    `---- static/
        |---- css/style.css          # Stile globale, layout, animazioni
        |---- js/                    # toast.js, lateral-toggle.js
        `---- images/                # logo
\end{verbatim}

\clearpage

\section{Funzionalit\`{a} implementate}

\subsection{Autenticazione}
Il proprietario accede al sistema tramite email e password, verificate con una
query SQL diretta sulla tabella \texttt{PROPRIETARIO}. La password \`{e}
confrontata in chiaro, scelta effettuata per un database locale e semplicità di implementazione. \\

La sessione viene mantenuta tramite il framework di sessioni di
Django (cookie), configurato per scadere alla chiusura del browser. Ogni form
\`{e} protetto da CSRF tramite il token integrato di Django.

\subsection{Interfaccia e navigazione}
La dashboard espone le tre aree applicative (Spiaggia, Ordini, Turni) tramite
pulsanti che caricano la relativa pagina via HTMX. \\
All'interno di ciascuna area, un menu laterale distingue le operazioni di inserimento/modifica
(``Servizi'') dalle query di controllo (``Controlli''): le prime vengono
mostrate in un pannello centrale, le seconde in un pannello laterale che
appare solo quando contiene un risultato. Un riclic sullo stesso pulsante
richiude il pannello corrispondente, riportando l'interfaccia allo stato
iniziale.

\subsection{Operazioni implementate}
Sono state implementate le stesse operazioni descritte nella sezione~\ref{sec:operazioniPrincipali}, 
con qualche aggiunta effettuata per comodità come:
\begin{enumerate}
  \item Check del \entita{Catalogo} di dato un \entita{Fornitore}
  \item Modifica di inserimento ordine e nuovo fornitore per aggiungere più prodotti alla volta 
  (l'aggiunta di un solo elemento a richiesta risultava ripetitivo)
\end{enumerate} 

% -- 5.6 Screenshot ------------------------------------------
\section{Screenshot dell'interfaccia}

\begin{figure}[H]
  \centering
  % \includegraphics[width=0.9\textwidth]{immagini/screen_login.png}
  \fbox{\parbox{0.9\textwidth}{\centering\vspace{3cm}
    \textit{[Screenshot: pagina di login]}
  \vspace{3cm}}}
  \caption{Pagina di login}
  \label{fig:screen-login}
\end{figure}

\begin{figure}[H]
  \centering
  % \includegraphics[width=0.9\textwidth]{immagini/screen_dashboard.png}
  \fbox{\parbox{0.9\textwidth}{\centering\vspace{3cm}
    \textit{[Screenshot: dashboard principale]}
  \vspace{3cm}}}
  \caption{Dashboard --- accesso alle tre aree applicative}
  \label{fig:screen-dashboard}
\end{figure}

\begin{figure}[H]
  \centering
  % \includegraphics[width=0.9\textwidth]{immagini/screen_spiaggia.png}
  \fbox{\parbox{0.9\textwidth}{\centering\vspace{3cm}
    \textit{[Screenshot: area spiaggia]}
  \vspace{3cm}}}
  \caption{Area Spiaggia --- pannello operazioni e risultati query}
  \label{fig:screen-spiaggia}
\end{figure}

\begin{figure}[H]
  \centering
  % \includegraphics[width=0.9\textwidth]{immagini/screen_turni.png}
  \fbox{\parbox{0.9\textwidth}{\centering\vspace{3cm}
    \textit{[Screenshot: gestione turni]}
  \vspace{3cm}}}
  \caption{Area Turni --- dipendenti e assegnazione turni}
  \label{fig:screen-turni}
\end{figure}

\begin{figure}[H]
  \centering
  % \includegraphics[width=0.9\textwidth]{immagini/screen_magazzino.png}
  \fbox{\parbox{0.9\textwidth}{\centering\vspace{3cm}
    \textit{[Screenshot: magazzino]}
  \vspace{3cm}}}
  \caption{Area Ordini --- controllo prodotti sotto scorta}
  \label{fig:screen-magazzino}
\end{figure}
